\section*{Cr\'edits, provenance et droits}
\begingroup\small
\noindent\textbf{Ouvrage original.} \textit{Cohomologie $\ell$-adique et fonctions $L$. S\'eminaire de G\'eom\'etrie Alg\'ebrique du Bois-Marie 1965--66 (SGA~5)}, dirig\'e par A.~Grothendieck, avec la collaboration de I.~Bucur, C.~Houzel, L.~Illusie, J.-P.~Jouanolou et J.-P.~Serre; \'edit\'e par L.~Illusie. Lecture Notes in Mathematics~589, Springer-Verlag, Berlin, Heidelberg, New York, 1977.

Les textes publi\'es repr\'esent\'es ici sont les expos\'es I, III, III~B, V, VI, VII, VIII, X, XII et XV, suivis de l'Index terminologique et de l'Index des notations. Les expos\'es I et III sont de A.~Grothendieck, r\'edig\'es par L.~Illusie; l'appendice \`a I et l'expos\'e III~B sont de L.~Illusie; les expos\'es V, VI et VII sont de J.-P.~Jouanolou; les expos\'es VIII, X et XII sont de A.~Grothendieck, r\'edig\'es par I.~Bucur; l'expos\'e XV est de C.~Houzel. Les attributions et r\'ef\'erences individuelles sont conserv\'ees dans le corps du texte.

\noindent\textbf{La pr\'esente \'edition.} Transcription fran\c{c}aise produite par IA (\textit{AI-transcribed}); traduction anglaise produite par IA (\textit{AI-translated}); corrections produites par IA (\textit{AI-produced corrections}) et composition produite par IA. Les sources fran\c{c}aises de r\'ef\'erence sont les deux fichiers d\'efinitifs \texttt{fr/sga5/diplomatic/s5.tex} et \texttt{fr/sga5/current/s5.tex} de \texttt{SGA5\_FR.zip}. Les documents du corpus Codex d'ao\^ut 2026 d\'efinissent les couches diplomatique et corrig\'ee et conservent les registres de corrections et de r\'evisions. Le contr\^ole ult\'erieur sur le fac-simil\'e, la reprise manuelle et la passe des conventions sont document\'es dans \texttt{fixes/}, \texttt{reports/} et \texttt{conventions/}. L'identit\'e du fac-simil\'e original est conserv\'ee dans ces documents: 496 pages physiques, 62\,025\,563 octets, empreinte SHA-256 consign\'ee dans les registres.

La base anglaise est \texttt{en/ported\_source/sources/native/SGA5/} dans \texttt{SGA5\_EN.zip}, apr\`es la fusion et le report des corrections fran\c{c}aises consign\'es au commit \texttt{d11a0c6} de sga-en. Le rapport de report cr\'edite Claude Opus~5.5 pour les scripts, la fusion et le contr\^ole manuel, et GPT-6.1 Sol, effort high, pour les d\'ecisions et les modifications propos\'ees. Les rapports fran\c{c}ais de reprise manuelle et de conventions cr\'editent Claude Opus~5.5. Les documents du d\'ep\^ot anglais cr\'editent OpenAI Codex, GPT-5.6 Sol pour la traduction et GPT-6 Astra pour les corrections et l'assemblage en septembre--octobre 2026.

\noindent\textbf{Pr\'eparation des pr\'esents lecteurs.} GPT-6 Astra (OpenAI, ChatGPT Pro mode). Les six lecteurs sont engendr\'es \`a partir d'un fichier ma\^itre annot\'e par langue. Les registres joints identifient la source de chaque variante et conservent les d\'ecisions relatives aux lectures encore discut\'ees.

\noindent\textbf{Contributions humaines conserv\'ees des documents de la collection.} Les attributions suivantes reprennent les contributions et leurs volumes respectifs consign\'es dans \texttt{en/repo/README.md} et \texttt{en/repo/ATTRIBUTION.md}. La collection originale comprend des travaux de Alexander Grothendieck, Mich\`ele Raynaud, Michel Raynaud, Pierre Deligne, Nicholas Katz, Michel Demazure, Michael Artin, Jean-Louis Verdier, Luc Illusie et leurs collaborateurs.

SGA~1 conserve la recomposition fran\c{c}aise associ\'ee au projet de b\'en\'evoles de Bas Edixhoven, \`a la Soci\'et\'e math\'ematique de France et aux notes de Michel Raynaud (arXiv \texttt{math/0206203v2}). La traduction anglaise de SGA~2 provient du TeX fran\c{c}ais corrig\'e de la r\'e\'edition dirig\'ee par Yves Laszlo, \textit{Cohomologie locale des faisceaux coh\'erents et th\'eor\`emes de Lefschetz locaux et globaux}, Documents Math\'ematiques~4, Soci\'et\'e math\'ematique de France, 2005 (arXiv \texttt{math/0511279v1}). La plus grande partie de la composition fran\c{c}aise initiale a \'et\'e r\'ealis\'ee par les b\'en\'evoles du projet de Bas Edixhoven: L.~Bayle, N.~Borne, O.~Brinon, J.~Buresi, M.~Chardin, F.~Ducrot, P.~Graftiaux, F.~Han, P.~Karwasz, L.~Koelblen, D.~Madore, S.~Morel, D.~Naie, B.~Osserman, J.~Riou et V.~S\'echerre; C.~Sabbah l'a adapt\'ee au format de la SMF.

Les sources anglaises de SGA~2 et SGA~3 reconnaissent la traduction anglaise de Jacob C.~Reinhold, instantan\'e \texttt{e7a259f3f8608ad3edf9bf6eead3fd504dd2d23e}, utilis\'ee pour la comparaison et la r\'edaction. Sa mention CC~BY~4.0 s'applique \`a cette contribution et est conserv\'ee; la formulation anglaise qui en provient demeure attribu\'ee \`a son auteur humain. Les lemmes, d\'emonstrations et corrections ajout\'es dans la r\'e\'edition de SGA~3 par Philippe Gille et Patrick Polo sont cr\'edit\'es \`a cette r\'e\'edition dans les notes et additions de la traduction.

SGA~4 est traduit \`a partir de l'\'edition fran\c{c}aise en LaTeX diffus\'ee par Fabrice Orgogozo, instantan\'e au commit \texttt{71766d9}, et emploie ses fichiers de configuration TeX. Les notes intitul\'ees \og N.D.E.\fg{} sont les notes de cette \'edition, traduites. SGA~$4\frac12$ conserve la composition LaTeX de Daniel Miller et sa note originale de compositeur. SGA~7~I suit le volume publi\'e et un PDF recompos\'e; \texttt{en/repo/SOURCE\_AUTHORITY.md} en consigne la num\'erotation des expos\'es. Le d\'ep\^ot cr\'edite aussi Claude Opus~5.5 (Anthropic) pour les sources de SGA~2 et SGA~3 du 26 octobre 2026 et la traduction de SGA~3 revue sur l'impression de 1970 en novembre 2026, ainsi que GPT-6.1 Sol (OpenAI API) pour les modifications propos\'ees reportant dans SGA~4 les corrections du contr\^ole fran\c{c}ais. Les attributions humaines propres \`a chaque source et les mentions authentiques de droits des tiers demeurent applicables.

\noindent\textbf{Droits.} Les droits sur l'ouvrage original appartiennent \`a son auteur ou aux repr\'esentants ou h\'eritiers de l'auteur. Dans la mesure o\`u nous d\'etenons des droits sur la transcription, la traduction, les corrections et la composition produites par IA, ces droits sont d\'edi\'es au domaine public sous CC0 1.0. Cette \'edition est publi\'ee \`a des fins \'educatives et non lucratives, et nous ne revendiquons aucun droit que nous ne d\'etenons pas. Si vous d\'etenez des droits sur cet ouvrage et vous opposez \`a sa publication ici, veuillez ouvrir une issue GitHub et nous le retirerons.

\noindent\textbf{Pages liminaires de l'imprim\'e.} La page de titre originale, l'Introduction de Luc Illusie (p.~iii--viii), la Liste des SGA et la Table des mati\`eres n'ont pas de transcription dans les sources fournies. La page de titre et la table des mati\`eres de cette \'edition sont cr\'e\'ees \`a partir du texte disponible; l'Introduction de l'expos\'e~I appartient \`a cet expos\'e.
\endgroup
\clearpage
